% Suffixes used in medical terminology.

\medterm{-Stomy} A suffix meaning surgical creation of an opening, such as colostomy or gastrostomy.
\medterm{-Tomy} A suffix meaning incision or cutting into.
\medterm{Cidal} A suffix meaning killing or destroying, as in bactericidal or fungicidal.

\medterm{Cide} A suffix meaning an agent or process that kills, as in pesticide, bactericide, or homicide.

\medterm{Cyte} A suffix meaning cell, as in erythrocyte or leukocyte.

\medterm{Cytosis} A suffix or term indicating a condition involving cells or, in some contexts, an increase in the number of a particular type of cell.

\medterm{Dactyl-, -Dactyl} A combining form or suffix meaning finger or toe.

\medterm{Ectomy} A suffix meaning surgical removal or excision, as in appendectomy or mastectomy.

\medterm{Emia} A suffix meaning a condition of the blood or a substance present in the blood, as in anemia or bacteremia.

\medterm{Iatrics} A suffix or word element referring to medical treatment or a medical specialty, as in geriatrics and pediatrics.

\medterm{Iatry} A suffix meaning medical treatment or a medical specialty, as in psychiatry and podiatry.

\medterm{Itis} A suffix meaning inflammation, as in dermatitis, bronchitis, and arthritis.

\medterm{Kalemia} A suffix or term referring to the concentration of potassium in the blood, as in hypokalemia or hyperkalemia.

\medterm{MAB} Abbreviation for monoclonal antibody; the suffix \textit{-mab} in a generic drug name indicates a monoclonal antibody, such as adalimumab or rituximab.

\medterm{Oma} -oma is a suffix commonly indicating a tumor or mass, although some words ending in -oma describe non-neoplastic conditions.

\medterm{Pathy} A suffix meaning disease, disorder, or abnormality, as in neuropathy, cardiomyopathy, and nephropathy.

\medterm{Plegia} -plegia is a suffix meaning paralysis, as in paraplegia, hemiplegia, or quadriplegia.

